\section{相关工作}

本节不按单篇论文罗列文献，而按论文故事所需的功能分组：先说明跨平台基础方法，再说明不同平台上的防晃研究，最后聚焦移动底盘液体运输和普通移动机器人局部规划器的差异。该结构服务于一个核心定位：\SPMPC{} 位于移动底盘液体运输中的在线滚动防晃局部规划方向。

\subsection{跨平台基础：液体模型、输入整形与测量评价}

液体防晃规划首先需要可用于控制和优化的液体模型。高保真 CFD、VOF、FEM 或 SPH 方法能够描述复杂自由液面、冲击和非线性现象，但其计算代价通常不适合直接嵌入实时局部规划器。因此，许多面向控制和规划的研究采用低阶模型，例如等效摆模型、质量-弹簧-阻尼模型或模态状态空间模型。这类模型的共同作用是用有限维状态表示液体对历史运动激励的动态记忆。

输入整形、命令预整形、速度剖面设计和 jerk 限制等方法是另一类重要基础。它们说明通过改变输入频率成分、速度变化和加速度形状可以减少残余振动或液面晃动。然而，这类方法通常依赖已知频率、固定路径或预先设计的 profile，并不必然传播当前液体状态。因此，它们适合作为 smooth/profile 类 baseline 的理论基础，但不能直接等价于在线 slosh-aware 局部规划。

此外，真实液面评价是防晃论文可信度的重要组成部分。若仅使用规划器内部的模型输出作为评价指标，容易产生“用自己的模型评价自己”的问题。因此，本文将区分内部预测的 model proxy 与外部真实液面评价，并优先采用 RGB 液面观测、max-LCR 或其他视觉/传感指标衡量真实晃动。

\subsection{机械臂、SCARA 与操作机器人中的液体防晃}

机械臂和 SCARA 等操作机器人中的液体防晃研究通常关注末端轨迹规划、容器姿态补偿、平坦性规划、预测控制或鲁棒输出反馈控制。此类方法表明，在机器人运动规划和控制中显式考虑液体响应是有价值的，也说明防晃不应只依赖末端轨迹的几何平滑性。

但是，机械臂类平台与移动底盘平台存在关键差异。机械臂通常能够控制末端姿态，甚至利用冗余自由度改变容器倾角，以主动抵消液面激励；而差速或 skid-steer 移动底盘主要通过线速度、角速度、加速度和角加速度影响容器运动。二者的控制输入、任务形式和动力学约束不同。因此，机械臂/SCARA 防晃文献在本文中属于 broader related work，而不是 Scout Mini 移动底盘的同层实验 baseline。

\subsection{液罐车、车辆与船舶等载液平台}

液罐车、车辆和船舶等大尺度载液平台中的晃动研究关注制动、转弯、侧倾稳定性、结构载荷和车辆-液体耦合动力学。相关研究说明，移动载体中的液体晃动会显著影响运动安全和系统稳定性，也常采用 preview control、悬架控制、有限元分析或流固耦合模型分析晃动影响。

这类工作为本文提供了移动载液系统的安全背景，但其问题尺度和控制目标通常不同于服务机器人或实验室移动底盘。液罐车/船舶更强调大尺度载荷、结构响应或车辆稳定性，而本文关注的是给定安全路径上的在线局部轨迹规划。因此，该类文献可作为 broader background，不作为本文同层方法对比的主线。

\subsection{移动底盘液体运输与防晃轨迹规划}

与本文最接近的是移动底盘或服务机器人液体运输中的防晃路径与轨迹规划研究。已有方法大致可以分为四类。

\begin{table}[t]
  \centering
  \caption{移动底盘液体运输近邻工作的层级划分}
  \label{tab:mobile_liquid_related_work}
  \begin{tabular}{p{0.18\linewidth} p{0.34\linewidth} p{0.36\linewidth}}
    \toprule
    类别 & 典型问题 & 与本文关系 \\
    \midrule
    固定路径/曲线路径设计 & 预先设计路径几何、转弯半径和曲率以减少晃动 & 说明路径几何会影响晃动，但多为预设计路径 \\
    速度剖面/input shaping & 在给定路径上设计速度、加速度或命令输入 & 支撑 smooth/profile baseline，但通常不是在线状态感知规划 \\
    离线防晃轨迹优化 & 将液体模型或约束放入整段轨迹优化 & 显式考虑液体，强近邻，但常偏离线或预规划 \\
    防晃跟踪控制/特殊机构 & 给定轨迹后进行 tracking 或利用主动抑振机构 & 说明防晃控制有效，但层级或硬件条件不同 \\
    \bottomrule
  \end{tabular}
\end{table}

这些方法为移动底盘液体运输提供了重要基础，但它们通常不是导航栈中的在线局部轨迹规划问题。本文希望强调的 gap 是：移动底盘在运行过程中需要不断更新下一段局部运动，并且此时液体当前相位和历史激励会影响未来的晃动响应。若规划器不传播液体状态，则无法在局部决策中利用这些动态信息。

\subsection{不含液体动态的移动机器人局部规划器}

普通移动机器人局部规划器和轨迹优化方法，如 DWA、TEB、pure pursuit、MPPI、通用轨迹优化、jerk-limited trajectory generation 和 time-optimal path parameterization，已经能够处理路径跟踪、运动可行性、平滑性甚至避障或风险约束。它们是本文实验中必须考虑的外部 baseline，因为它们代表了“如果不显式考虑液体，普通在线 planner 能做到什么”。

然而，这些方法的目标函数和状态变量通常不包含被运输液体的模态状态。即使它们生成的轨迹速度和加速度较平滑，也无法表示液体的当前相位、残余振荡和对未来输入的动态响应。因此，它们不是液体防晃方向本身，而是用于验证本文 claim 的 negative control：ordinary online planning is not necessarily slosh-aware planning。

\subsection{SPMPC 的位置}

综上，现有防晃研究证明了液体模型、输入整形、轨迹优化和防晃控制的价值；机械臂和载液车辆平台说明液体晃动会影响机器人和移动载体任务；移动底盘液体运输近邻方法则覆盖了固定路径、速度剖面、离线优化和跟踪控制等方向。但是，面向移动底盘开口液体运输的在线滚动防晃局部规划仍是一个需要明确表述的 gap。

本文的 \SPMPC{} 位于该 gap：它不是普通 local planner 加一个平滑代价，也不是给定轨迹后的控制器，而是在移动底盘局部轨迹规划阶段显式传播低阶液体模态状态，并在滚动时域内联合优化路径进度、底盘控制和预测液体响应。
